% 
%
% (c) 2025 Autor, ETH Zürich
%
% !TEX root = linalg_I-II_summary.tex
% !TEX encoding = UTF-8
%

\section{Determinants}
Let 
\[
	A = \begin{pmatrix}
		a & b\\
		c & d
	\end{pmatrix} \in M_{2 \times 2}(K).
\]
There is a very simple criterion to check wether $A$ is invertible or not.

\begin{proposition}{Criterion for Invertibility of a $2 \times 2$ Matrix}{crit_inv_2by2_mat}
	$A$ is invertible if and only if $ad - bc \neq 0$. Moreover, if $A$ is invertible, then
	\[
		A^{-1} = \frac{1}{ad - bc}\, \begin{pmatrix}
			d & -b\\
			-c & a
		\end{pmatrix}.
	\]
	The number $ad - bc$ is $\det(A)$.
\end{proposition}

\begin{definition}{N-Linear Function}{nlin_func}
	Let $D:M_{n\times n}(K) \to K$ be a function. We say that $D$ is \textbf{$n$-linear} if for every $1 \leq i \leq n$, $D$ is a linear function of the $i$-th row when the other $(n-1)$ rows are held fixed.
\end{definition}
We can think of $M_{n\times n }(K)$ as $K_{\text{row}}^n \times \hdots \times K_{\text{row}}^n$ and if
\[
	A = \begin{pmatrix}
		\rule[0.5ex]{1em}{0.4pt} & \alpha_1 & \rule[0.5ex]{1em}{0.4pt}\\
		& \vdots & \\
		\rule[0.5ex]{1em}{0.4pt} & \alpha_n & \rule[0.5ex]{1em}{0.4pt}
	\end{pmatrix}
\]
we write $D(A) = D(\alpha_1, \hdots , \alpha_n)$. Saying that $D$ is $n$-linear means that for all $1 \leq i \leq n$, we have that
\begin{enumerate}
	\item \begin{align*}
		D&(\alpha_1, \hdots , \alpha_i + \alpha_i', \hdots , \alpha_n) = D(\alpha_1, \hdots , \alpha_i, \hdots , \alpha_n) + D(\alpha_1, \hdots , \alpha_i', \hdots , \alpha_n) \\
		\forall& \, \alpha_1, \hdots , \alpha_i, \alpha_i', \hdots , \alpha_n \in K^n_{\text{row}}
	\end{align*}
	\item \begin{align*}
		D&(\alpha_1, \hdots , c \cdot \alpha_i, \hdots, \alpha_n) = c \cdot D(\alpha_1, \hdots , \alpha_i, \hdots, \alpha_n) \\
		\forall& \,c \in K, \,\alpha_1, \hdots \alpha_i, \hdots , \alpha_n \in K^n_{\text{row}}
	\end{align*}
\end{enumerate}

\begin{lemma}{Linear Combinations of $n$-linear Functions}{lin_comb_nlin_func}
	Any linear combination of $n$-linear functions is an $n$-linear function again.
\end{lemma}

\begin{definition}{Alternating Functions}{alt_func}
	Let $D:M_{n\times n}(K) \to K$ be an $n$-linear function. We say that $D$ is \textbf{alternating} if for every matrix $A \in M_{n\times n}(K)$ in which some two rows are equal we have that $D(A) = 0$.
\end{definition}

\begin{proposition}{}{prop_alt}
	Let $D$ be an $n$-linear function. Assume that $D$ has the property that whenever $A$ has two equal adjacent (aufeinanderfolgende) rows then $D(A) = 0$. Then
	\begin{enumerate}
		\item $D$ is alternating.
		\item For every matrix $B$, if $B'$ is obtained from $B$ by interchanging any two of its rows, then $D(B') = -D(B)$.
	\end{enumerate}
\end{proposition}

\begin{definition}{Determinant Function}{det_func}
	A function $D:M_{n\times n}(K) \to K$ is called \textbf{determinant function} if $D$ is $n$-linear, is alternating and $D(I) = 1$.
\end{definition}

\begin{definition}{Row\&Column-Deletion Matrix}{row_col_del_mat}
	Let $A \in M_{n\times n}(K)$, $n \geq 2$, $1 \leq i \leq n$, $1 \leq j \leq n$. We denote by 
	\[
		A(i|j) \in M_{(n-1)\times (n-1)}(K)
	\]
	the matrix obtained from $A$ by deleting row \#$i$ and column \#$j$.
	
	If $D$ is an $(n-1)$-linear function and $A \in M_{n\times n}(K)$ we denote
	\[
		D_{ij}(A) := D\big(A(i|j)\big).
	\]
\end{definition}

\textbf{Notation:} Let $A \in M_{m\times n}(K)$ and $1 \leq i \leq m$, $1 \leq j \leq n$. Denote the element at position $(i,j)$ in the matrix $A$ by $A(i,j)$ or $A_{ij}$.

\begin{theorem}{Entwicklungssatz}{entw_satz}
	Let $n \geq 2$ and $D$ an $(n-1)$-linear function, which is alternating. Let $1 \leq j \leq n$. Define a function $E_j:M_{n\times n}(K) \to K$ by
	\[
		E_j(A) := \sum_{i=1}^{n} (-1)^{i+j} A_{ij} \, D_{ij}(A).
	\]
	Then $E_j$ is an $n$-linear function and is alternating. If $D$ is a determinant function then so is $E_j$.
\end{theorem}

\begin{corollary}{Existence of a Determinant Function}{exist_det_func}
	$\forall \, n \geq 1$, there exists at least one determinant function $M_{n\times n}(K) \to K$.
\end{corollary}

\subsection{Uniqueness of the Determinant}
Let $D:M_{n\times n}(K) \to K$ be an $n$-linear alternating function. Let
\[
	A = \begin{pmatrix}
		\rule[0.5ex]{1em}{0.4pt} & \alpha_1 & \rule[0.5ex]{1em}{0.4pt}\\
		& \vdots & \\
		\rule[0.5ex]{1em}{0.4pt} & \alpha_n & \rule[0.5ex]{1em}{0.4pt}
	\end{pmatrix} \in M_{n\times n}(K).
\]
Denote
\[
	I = \begin{pmatrix}
		1 & \hdots & 0\\
		\vdots &\ddots & \vdots\\
		0 & \hdots & 1
	\end{pmatrix}
	=
	\begin{pmatrix}
		\rule[0.5ex]{1em}{0.4pt} & \varepsilon_1 & \rule[0.5ex]{1em}{0.4pt}\\
		& \vdots & \\
		\rule[0.5ex]{1em}{0.4pt} & \varepsilon_n & \rule[0.5ex]{1em}{0.4pt}
	\end{pmatrix}.
\]
(So, $\varepsilon_1, \hdots , \varepsilon_n$ is the standard basis of $K^n_{\text{row}}$.) Clearly, 
\[
	\alpha_i = \sum_{j=1}^{n} A(i,j)\, \varepsilon_j \qquad \forall \, 1\leq i \leq n.
\]
Therefore,
\begin{align*}
	D(A) &= D(\alpha_1, \hdots , \alpha_n) = D\bigg( \sum_{j=1}^{n} A(1,j) \, \varepsilon_j, \alpha_2, \hdots , \alpha_n\bigg) = \sum_{j=1}^{n} A(1,j) \cdot D(\varepsilon_j, \alpha_2, \hdots , \alpha_n) \\
	&= \sum_{j=1}^{n} A(1, j)\cdot D\bigg(\varepsilon_j, \sum_{k=1}^{n} A(2, k)\, \varepsilon_k, \hdots, \alpha_n\bigg) = \sum_{j,k = 1}^{n} A(1,j)\cdot A(2,k) \cdot D(\varepsilon_j, \varepsilon_k, \hdots, \alpha_n).
\end{align*}
Applying these steps to every row we get that
\begin{equation}
	\label{eq:det_func}
	D(A) = \sum_{k_1, \hdots, k_n = 1}^{n} A(1, k_1)\cdot A(2, k_2)\cdot \hdots \cdot A(n, k_n)\cdot D(\varepsilon_{k_1}, \varepsilon_{k_2} \hdots, \varepsilon_{k_n}),
\end{equation}
where each of the indices $k_1, \hdots , k_n$ runs between $1$ and $n$. So far we haven't used the fact that $D$ is alternating. Now, since $D$ is alternating $D(\varepsilon_{k_1}, \hdots , \varepsilon_{k_n}) = 0$, whenever two of the indices $k_1, \hdots , k_n$ are equal. So, we are interested only in ordered sequences of indices $(k_1, \hdots, k_n)$ in which $k_i \in \{1, \hdots, n\} \quad \forall i$ and $k_i \neq k_j \quad \forall i \neq j$. Such a sequence is called a permutation of $\{1, \hdots , n\}$, or a permutation of degree $n$.

\begin{definition}{Permutation}{permut}
	A \textbf{permutation} is a bijective map $\sigma :\{1, \hdots, n\} \to \{1, \hdots, n\}$.
\end{definition}

So Alternatively, we can think of a permutation of degree $n$ as a bijective function $\sigma:\{1,\hdots , n\} \to \{1, \hdots , n\}$. Then the sequence $(k_1, \hdots, k_n) = (\sigma 1, \hdots, \sigma 2)$ is a permutation as defined earlier. (Notation: $\sigma(i) := \sigma i$.). In other words, a permutation $\sigma$ of degree $n$ is a re-ordering of a list with $n$ elements $(1,\hdots, n) \overset{\sigma}{\longrightarrow} (\sigma 1, \hdots, \sigma n)$.

We can rewrite Equation \eqref{eq:det_func} as
\[
	D(A) = \sum_{\sigma \in S_n} A(1, \sigma 1)\cdot \hdots \cdot A(n, \sigma n) \cdot D(\varepsilon_{\sigma 1}, \hdots , \varepsilon_{\sigma n}),
\]
where $S_n$ is the set of all permutations of degree $n$.

\subsubsection{Important Thing About Permutations}
Let $\sigma$ be a permutation of degree $n$. Then one can pass from $(1, \hdots , n)$ to $(\sigma 1, \hdots , \sigma n)$ by performing a finite sequence of interchanges of pairs, called transpositions.

\begin{definition}{Transposition}{transposition}
	A \textbf{transposition} is a permutation $\theta :\{1, \hdots, n\} \to \{1, \hdots, n\}$, where $\exists \, i,j \in \{1, \hdots, n\}$, $i \neq j$ s.t. $\theta(i) = j$,  $\theta(j) = i$ and $\theta(k) = k \quad \forall \, k\neq i,j$.
\end{definition}

Now, every permutation $\sigma :\{1, \hdots, n\} \to \{1, \hdots, n\}$ can be written as a composition of transpositions
\begin{equation}
	\label{eq:permut_transp}
	\sigma = \theta_1 \circ \hdots \circ \theta_m,
\end{equation}
where $\theta_1, \hdots, \theta_m$ are transpositions. But these $\theta_i$'s and $m$ are not unique. In other words, the same $\sigma$ can be written as in Equation \eqref{eq:permut_transp} in many different ways
\[
	\sigma = \theta_1 \circ \hdots \circ \theta_m = \theta_1' \circ \hdots \circ \theta_{m'}'.
\]
Why Can $\sigma$ be written as in \eqref{eq:permut_transp}. Start with $(1, \hdots , n)$ and if $\sigma 1 \neq 1$ then interchange $1$ and $\sigma 1$. We now have $(\sigma 1, \hdots , 1, \hdots , n)$. If $\sigma 2 \neq 2$, then interchange $1 $ and $\sigma 2$; otherwise proceed to $3$, etc. But of course there are other ways to do it, e.g., $(3, 2, 1)$ can be obtained in 3 steps: $(1,2,3) \to (2,1,3) \to (2,3,1) \to (3,2,1)$. But also in one step $(1,2,3) \to (3,2,1)$.

\begin{definition}{\textcolor{red}{Sign of a Permutation}}{sgn_permut}
	Let $\sigma$ be a permutation and let $\theta_1, \hdots, \theta_m$ be transpositions s.t. $\sigma = \theta_1 \circ \hdots \circ \theta_m$. The permutation $\sigma$ is called \textbf{even} if $m$ is even, and is called \textbf{odd} if $m$ is odd. The number $(-1)^m \in \{-1, 1\}$ is called the \textbf{sign} of $\sigma$, i.e., $\operatorname{sgn}(\sigma) := (-1)^m$.
\end{definition}

\begin{lemma}{\textcolor{red}{Well Definedness of the Sign of a Permutation}}{parity_permut}
	If $(\sigma 1, \hdots , \sigma n)$ is obtained from $(1, \hdots , n)$ in two different ways, one time by a sequence of $m$ transpositions, and another time by a sequence of $m'$ transpositions, then $m'$ and $m$ have the same parity, i.e., $(-1)^{m'} = (-1)^m$.
	
	(So the parity of $m$ depends only on $\sigma$ and not on the particular way we performed a sequence of transpositions $(1, \hdots , n) \to \hdots \to (\sigma 1, \hdots, \sigma n)$.)
\end{lemma}

\begin{proof}
	We know that determinant functions $M_{n\times n}(K) \to K$ exists for every $n \geq 1$. Fix a determinant function $E:M_{n\times n}(K) \to K$. Consider $E(\varepsilon_{\sigma 1}, \hdots , \varepsilon_{\sigma n})$. If $(\sigma 1, \hdots, \sigma n)$ can be obtained from $(1, \hdots, n)$ by a sequence of $m$ transpositions, then the matrix
	\[
		\begin{pmatrix}
			\rule[0.5ex]{1em}{0.4pt} & \varepsilon_1 & \rule[0.5ex]{1em}{0.4pt}\\
			& \vdots & \\
			\rule[0.5ex]{1em}{0.4pt} & \varepsilon_n & \rule[0.5ex]{1em}{0.4pt}
		\end{pmatrix}
	\]
	can be obtained from
	\[
		I = \begin{pmatrix}
			1 & \hdots & 0\\
			\vdots &\ddots & \vdots\\
			0 & \hdots & 1
		\end{pmatrix}
	\]
	by a sequence of $m$ interchanges of pairs of rows. Therefore,
	\[
		E(\varepsilon_{\sigma 1}, \hdots, \varepsilon_{\sigma n}) = (-1)^m E(\varepsilon_1, \hdots, \varepsilon_n) = (-1)^m E(I) = (-1)^m.
	\]
	But this applies also to $m'$. So, 
	\[
		E(\varepsilon_{\sigma 1}, \hdots, \varepsilon_{\sigma n}) = (-1)^{m'}E(\varepsilon_1, \hdots, \varepsilon_n) = (-1)^{m'} E(I) = (-1)^{m'}.
	\]
	So we have that
	\[
		(-1)^m = (-1)^{m'}.
	\]
	\qedhere
\end{proof}

\subsubsection{Back to our $D(A)$}
We have seen that if $D$ is $n$-linear and alternating then
\[
	D(A) = \sum_{\sigma \in S_n} A(1, \sigma 1)\cdot \hdots \cdot A(n, \sigma n) \cdot D(\varepsilon_{\sigma 1}, \hdots , \varepsilon_{\sigma n}).
\]
But 
\begin{equation}
	\label{eq:det_permut}
	D(\varepsilon_{\sigma 1}, \varepsilon_{\sigma n}) = (\operatorname{sgn}\sigma)\cdot D(\varepsilon_1, \hdots, \varepsilon_n) = (\operatorname{sgn}\sigma)\cdot D(I)
\end{equation}

\begin{theorem}{Uniqueness of the Determinant Function}{unique_det_func}
	$\forall \, n \geq 1$, there exists a unique determinant function $M_{n\times n}(K) \to K$. We will denote it by $\det$. It is given by the formula
	\[
		\det(A) = \sum_{\sigma \in S_n} A(1, \sigma 1)\cdot \hdots \cdot A(n, \sigma n).
	\]
\end{theorem}

\begin{corollary}{}{cor_nlin_alt_func}
	For every $n$-linear and alternating function $D:M_{n\times n}(K) \to K$ we have that
	\[
		D(A) = \det(A)\cdot D(I) \qquad \forall \, A \in M_{n\times n}(K).
	\]
\end{corollary}

\subsubsection{A Bit more on Permutations}
Denote by $S_n$ the set of permutations $\sigma :\{1, \hdots, n\} \to \{1, \hdots, n\}$. $S_n$ has the structure of a group:
\begin{itemize}
	\item[$\bullet$] We can compose permutations, i.e., if $\sigma, \tau :\{1, \hdots, n\} \to \{1, \hdots, n\}$, then $\sigma \cdot \tau := \sigma \circ \tau$, where $\sigma \circ \tau :\{1, \hdots, n\} \to \{1, \hdots, n\}$.
	\item[$\bullet$] The identity $id :\{1, \hdots, n\} \to \{1, \hdots, n\}$ is the neutral element $1 \in S_n$. 
	\item[$\bullet$] The inverse $\sigma^{-1}$ of $\sigma$, is given by the inverse of the function $\sigma$ (recall that $\sigma$ is bijective, so $\sigma^{-1}$ does exist).
\end{itemize}
Also, note that $|S_n| = n!$

\begin{lemma}{Sign of the Composition of Permutations}{sgn_comp_permut}
	$\forall \, \sigma, \tau \in S_n$ we have 
	\[
		\operatorname{sgn}(\sigma \cdot \tau) = \operatorname{sgn}(\sigma)\cdot \operatorname{sgn}(\tau).
	\]
\end{lemma}

\begin{theorem}{\textcolor{red}{The Determinant of the Product of Two Matrices}}{det_prod_mat}
	$\forall \, A,B \in M_{n\times n}(K)$ we have that
	\[
		\det(A\cdot B) = \det(A) \cdot \det(B).
	\]
\end{theorem}

\begin{proof}
	Fix a matrix $B \in M_{n\times n}(K)$, and consider a function $D:M_{n\times n}(K) \to K$ defined by
	\[
		D(A) := \det(A\cdot B).
	\]
	
	\textbf{Claim:} $D$ is $n$-linear and alternating.
	
	\textbf{Proof of Claim:} Write
	\[
		A = \begin{pmatrix}
			\rule[0.5ex]{1em}{0.4pt} & \alpha_1 & \rule[0.5ex]{1em}{0.4pt}\\
			& \vdots & \\
			\rule[0.5ex]{1em}{0.4pt} & \alpha_n & \rule[0.5ex]{1em}{0.4pt}
		\end{pmatrix}.
	\]
	If we view $\alpha_i$ as a $1 \times n$ matrix, then $\alpha_i B$ is also $1 \times n$, and in fact
	\[
		A\cdot B = \begin{pmatrix}
			\rule[0.5ex]{1em}{0.4pt} & \alpha_1 B & \rule[0.5ex]{1em}{0.4pt}\\
			& \vdots & \\
			\rule[0.5ex]{1em}{0.4pt} & \alpha_n B & \rule[0.5ex]{1em}{0.4pt}
		\end{pmatrix}.
	\]
	So, 
	\[
		D(A) = \det(\alpha_1 B, \hdots, \alpha_n B).
	\]
	Clearly, for every two row vectors $\alpha_i, \alpha_i'$ we have that $(\alpha_i + \alpha_i') B = \alpha_i B + \alpha_i' B$ and
	\[
		\forall \, c\in K : \; (c \cdot \alpha_i) B = c \cdot (\alpha_i B).
	\]
	Since $\det$ is $n$-linear it follows that $D$ is also $n$-linear.
	
	If $\alpha_i = \alpha_j$, then $\alpha_i B = \alpha_j B$. Therefore, 
	\[
		D(A) = \det(\alpha_1 B, \hdots, \alpha_n B) = 0.
	\]
	This proves the claim.
	
	We continue with the proof of the Theorem. Since we found that $D$ is $n$-inear and alternating, by Corollary \ref{cor*cor_nlin_alt_func}, we have that
	\[
		D(A) = \det(A)\cdot D(I).
	\]
	But $D(I) = \det(I\cdot B) = \det(B)$. Therefore, 
	\[
		D(A) = \det(A \cdot B) = \det(A)\cdot \det(B). \qedhere
	\]
\end{proof}

\begin{corollary}{Equivalent Condition for The Invertibility of a Matrix}{equiv_cond_inv_mat}
	If $A$ is invertible then $\det(A) \neq 0$. Moreover,
	\[
		\det(A^{-1}) = \frac{1}{\det(A)}.
	\]
\end{corollary}

\subsection{More on Determinants}

\begin{theorem}{Determinant of a Transposed Matrix}{det_trans_mat}
	$\forall \, A \in M_{n\times n}(K) : \; \det(A^T) = \det(A)$.
\end{theorem}

\begin{corollary}{}{cor_alt_func_2}
	If $D:M_{n\times n}(K) \to K$ is an alternating $n$-linear function (of the rows of a matrix) then $D$ is also an alternating and $n$-linear function of the columns of a matrix. In fact, $D(A^T) = D(A) \quad \forall \, A\in M_{n\times n}(K)$. In particular, this holds also for $D = \det$.
\end{corollary}

\begin{theorem}{Determinant after Row/Column Operations}{det_row_op}
	\begin{enumerate}
		\item If $B$ is obtained from $A$ by adding a multiple of one row of $A$ to another row (or by doing the analogous operation on columns), then $\det(B) = \det(A)$.
		\item If $A \overset{R_i \leftrightarrow R_j}{\longrightarrow} B$, $i \neq j$, (or the analogous operation on columns) then $\det(B) = -\det(A)$.
		\item If $A \overset{c\cdot R_i \to R_i}{\longrightarrow} B$, then $\det(B) = c \cdot \det(A)$.
	\end{enumerate}
\end{theorem}

\subsubsection{Determinants of Block Matrices}
Consider a block matrix $M$ of the type
\[
	M = \begin{pmatrix}
		A & B\\
		0 & C
	\end{pmatrix},
\]
where $A \in M_{r\times r}(K)$, $C \in M_{s \times s}(K)$, $B \in M_{r \times s}(K)$.

\begin{proposition}{Determinant of a Block Matrix}{det_block_mat}
	Let $M$ be as above. Then
	\[
		\det(M) = \det(A) \cdot \det(C).
	\]
\end{proposition}

\begin{corollary}{Entwicklungssatz with Cofactors}{entwick_satz_cofact}
	$\forall \, A \in M_{n\times n}(K)$ we have the following $n$ different formulas for $\det(A)$, i.e.,
	\[
		\det(A) = \sum_{i=1}^{n} (-1)^{i+j} A_{ij} \det\big(A(i|j)\big), \quad j = 1, \hdots, n.
	\]
	The scalars $c_{ij} := (-1)^{i+j} \det\big(A(i|j)\big)$ are called the $i, j$ \textbf{cofactor} of $A$. We have
	\begin{equation}
		\label{eq:det_cofactor}
		\det(A) = \sum_{i=1}^{n} A_{ij}\cdot c_{ij}.
	\end{equation}
\end{corollary}

\textbf{Claim:} If we replace $A_{ij}$ in Equation \eqref{eq:det_cofactor} by $A_{ik}$, where $k$ is fixed, we always get 0. In other words, $\forall \, A \in M_{n\times n}(K)$, $\forall \, j$ we have
\[
	\sum_{i=1}^{n} A_{ik} \cdot c_{ij} = \begin{cases}
		\det(A) & j = k\\
		0 & j \neq k
	\end{cases}.
\]

\textbf{Proof of Claim:} Let $B \in M_{n \times n}(K)$ be the matrix obtained from $A$ by replacing column $j$ of $A$ by column $k$ of $A$. So, in $B$ column $j$ equals column $k$. Clearly, $\det(B) = 0$ since $B$ has two equal columns ($\Longleftrightarrow \; B^T$ has two equal rows). 
\[
	\Longrightarrow \qquad 0 = \det(B) = \sum_{i=1}^{n} (-1)^{i+j} B_{ij} \det\big(B(i|j)\big)= \sum_{i=1}^{n}  = \sum_{i=1}^{n} A_{ik} \det\big(A(i|j)\big) = \sum_{i=1}^{n} A_{ik} \cdot c_{ij}.
\]
This proves the claim.

So, $\forall \, A \in M_{n\times n}(K)$, $\forall \, j,k$ we have 
\begin{equation}
	\label{eq:det_cofactor_2}
	\sum_{i=1}^{n} A_{ik} \cdot c_{ij} = \delta_{jk} \cdot \det(A).
\end{equation}

\begin{definition}{Classical Adjoint of $A$}{classic_adjoint}
	The matrix $(c_{ij})^T$ is called the \textbf{classical adjoint of $A$}, i.e.,
	\[
		\operatorname{adj}(A) \in M_{n \times n}(K), \qquad \big(\operatorname{adj}(A)\big)_{ij} = c_{ji} = (-1)^{i+j} \cdot \det\big(A(j|i)\big).
	\]
\end{definition}

The Equation \eqref{eq:det_cofactor_2} says that
\[
	\label{eq:adjoint}
	(\operatorname{adj}A) \cdot A = \det(A) \cdot I.
\]
Indeed, 
\begin{align*}
	 \det(A) \cdot \delta_{jk} &= \sum_{i=1}^{n} c_{ij} \cdot A_{ik} = \sum_{i=1}^{n} \big(\operatorname{adj}(A)\big)_{ji} \cdot A_{ik} \\
	 \Longrightarrow \qquad \det(A) \cdot I &= \operatorname{adj}(A) \cdot A.
\end{align*}

\begin{theorem}{Formula for the inverse of $A$}{formula_inv}
	Let $A \in M_{n\times n}(K)$. Then $A$ is invertible if and only if $\det(A) \neq 0$. When $A$ is invertible we have
	\[
		A^{-1} = \frac{1}{\det(A)} \, \operatorname{adj}(A).
	\]
\end{theorem}

\subsection{Cramer's Rule}
Let $A \in M_{n\times n}(K)$. We would like a formula to solve the system of equations $A x = b$, where
\[
	x = \begin{pmatrix}
		x_1\\
		\vdots \\
		x_n
	\end{pmatrix}
\]
are the unknowns and
\[
	b= \begin{pmatrix}
		b_1\\
		\vdots \\
		b_n
	\end{pmatrix} \in K^n
\]
is given. Assume $Ax = b$. Then
\begin{align*}
	\operatorname{adj}A \cdot A x &= \operatorname{adj}A \cdot b\\
	\Longrightarrow \qquad \det(A) \cdot x &= \operatorname{adj}\cdot b\\
	\Longrightarrow \qquad \forall \, 1 \leq j \leq n : \; \det(A)\cdot x_j &= \sum_{i=1}^{n} (\operatorname{adj}A)_{ji}\cdot  b_i = \sum_{i=1}^{n} (-1)^{i+j} b_i \det\big(A(i|j)\big),
\end{align*}
where the last expression is the determinant of the matrix obtained from $A$ by replecing column $j$ by $b$.

\begin{theorem}{Cramer's Rule}{cramers_rule}
	Let $A \in M_{n\times n}(K)$ with $\det(A) \neq 0$. Let $b \in K^n$. Then the system $A x = b$ has a unique solution, and it is given by the formula
	\[
		x_j = \frac{\det\big(A(j, b)\big)}{\det(A)} \qquad \forall \, 1 \leq j \leq n,
	\]
	where $A(j, b)$ is the matrix obtained from $A$ by replacing column $j$ by $b$.
\end{theorem}

\subsection{Determinants of Endomorphisms}
Let $A \in M_{n \times n}(K)$ and $P \in GL_n(K)$, then 
\[
	\det(P^{-1} A P) = \det(P^{-1}) \cdot \det(A) \cdot \det(P) = \frac{1}{\det(P)}\cdot \det(A) \cdot \det(P) = \det(A).
\]
In other words, similar matrices have the same determinant.

Let $V$ be a finite dimensional vector space over $K$. Put $n := \dim V$. Let $T:V \to V$ be a linear map. Pick a basis $\mathcal{B}$ for $V$. We have a matrix $[T]^{\mathcal{B}}_{\mathcal{B}} \in M_{n\times n}(K)$. Consider $\det [T]^{\mathcal{B}}_{\mathcal{B}} \in K$.

\textbf{Question:} Does $\det [T]^{\mathcal{B}}_{\mathcal{B}}$ depend on the basis $\mathcal{B}$?

\textbf{Answer:} Let $\mathcal{B}'$ be another basis for $V$. Denote by $P := [id_V]^{\mathcal{B}'}_{\mathcal{B}}$ the transition matrix between the basis $\mathcal{B}'$ and $\mathcal{B}$. Recall that
\[
	[T]^{\mathcal{B}'}_{\mathcal{B}} = [id_V]^{\mathcal{B}}_{\mathcal{B}'} \cdot [T]^{\mathcal{B}}_{\mathcal{B}} \cdot [id_V]^{\mathcal{B}'}_{\mathcal{B}},
\]
and we also have $[id_V]^{\mathcal{B}}_{\mathcal{B}'} = \big([id_V]^{\mathcal{B}'}_{\mathcal{B}}\big)^{-1}$. So, 
\[
	[T]^{\mathcal{B}'}_{\mathcal{B}'} = P^{-1} \cdot [T]^{\mathcal{B}}_{\mathcal{B}} \cdot P. \qquad \Longrightarrow \qquad \det [T]^{\mathcal{B}'}_{\mathcal{B}'} = \det [T]^{\mathcal{B}}_{\mathcal{B}}.
\]
Recall also that if $T, S:V \to V$ are two linear maps then $[S \circ T]^{\mathcal{B}}_{\mathcal{B}} = [S]^{\mathcal{B}}_{\mathcal{B}} \cdot [T]^{\mathcal{B}}_{\mathcal{B}}$. Recall also that if $T:V \to V$ is an isomorphism then $[T^{-1}]^{\mathcal{B}}_{\mathcal{B}} = \big([T]^{\mathcal{B}}_{\mathcal{B}}\big)^{-1}$. Finally, $[id_V]^{\mathcal{B}}_{\mathcal{B}} = I$.

\begin{corollary}{\textcolor{red}{Well Definedness of the Determinant of an Endomorphism}}{det_endo_well_def}
	Let $V$ be a finite dimensional vector space over $K$. Let $T:V \to V$ be a linear map. Then $\det [T]^{\mathcal{B}}_{\mathcal{B}}$ does \textbf{not} depend on the choice of basis $\mathcal{B}$ of $V$. We denote this scalar by $\det(T) \in K$. The function
	\[
		\det: \operatorname{End}(V) \to K, \quad T \mapsto \det(T)
	\]
	has the following properties:
	\begin{enumerate}
		\item $\det(S \circ T) = \det(S) \cdot \det(T) \qquad \forall \, S, T \in \operatorname{End}(V)$.
		\item $T$ is an isomorphism if and only if $\det(T) \neq 0$. If $T$ is an isomorphism then 
		\[
			\det(T^{-1}) = \frac{1}{\det(T)}.
		\]
		\item $\det(id_V) = 1$.
		\item $\det(0) = 0 \in K$.
	\end{enumerate}
\end{corollary}

\begin{remark}
	In contrast to $T:V \to V$, if we take $T:V \to W$, and $\mathcal{B}, \mathcal{C}$ bases of $V, W$, then $\det [T]^{\mathcal{B}}_{\mathcal{C}}$ \textbf{depends} on $\mathcal{B}, \mathcal{C}$ (unless $W = V$ and $\mathcal{C} = \mathcal{B}$).
\end{remark}